\section{Finite certificates}\label{sec:certificates}

This appendix completes the finite evidence used in
Lemma~\ref{interpolation:finite} and Lemma~\ref{signs:bernstein}.
We give the exact inputs, prove a uniform error bound for the finite
quadrature sums, and state the interval comparisons for the matrices,
residuals, and Bernstein coefficients.  The final subsection records the
scalar substitutions used for the infinite tails and continuum signs.
All terminating decimals denote exact rational numbers.

The primary calculation uses interval evaluation of the $N=256$ quadrature
sums, enlarged by the analytic moment remainder proved below.  It therefore
reduces the required numerical evidence to bounded finite sums, without
assuming convergence of a numerical integration or optimization procedure.
The infinite inverse and the signs between and beyond interpolation nodes
remain consequences of the analytic arguments in the main text.
Appendix~\ref{ind:section} gives the rational alternatives and their
separately stated scope.

\subsection{Exact input lists and matrices}\label{cert:inputs}

Table~\ref{tab:coefficients} gives the coefficients at every node of
$I_f=\mathcal N\cap[1,100]$.  Its four coefficient columns are multiplied
by $10^{10}$: an integer $a$ in the table means the coefficient
$a\,10^{-10}$.  At all nodes $n>100$ both approximate lists are zero.
The data at zero and the constant terms are, exactly,
\begin{equation}\label{cert:fixed-inputs}
 (\widetilde c_{1,0},\widetilde d_{1,0},C_1)
   =(1,.44,-.013),\qquad
 (\widetilde c_{2,0},\widetilde d_{2,0},C_2)
   =(0,-.368,.017).
\end{equation}
The zero-node data are included in all evaluations of the tabulated functions.
They are excluded when forming a column of the operator $S$.
Direct integer addition gives
\begin{equation}\label{cert:table-norms}
\begin{aligned}
 \sum_{n\ge0}(|\widetilde c_{1,n}|+|\widetilde d_{1,n}|)
   &=2.2516666247<2.29,\\
 \sum_{n\ge0}(|\widetilde c_{2,n}|+|\widetilde d_{2,n}|)
   &=.8927477314<2.29,\\
 \sum_{n>40}(|\widetilde c_{1,n}|+|\widetilde d_{1,n}|)
   &=.0000197471<.00002,\\
 \sum_{n>40}(|\widetilde c_{2,n}|+|\widetilde d_{2,n}|)
   &=.0000161623<.00002.
\end{aligned}
\end{equation}
Here and below an index $n\ge0$ in a coefficient sum ranges over
$\{0\}\cup\mathcal N$.

\begingroup
\small
\setlength{\tabcolsep}{7pt}
\begin{longtable}{r r r r r}
\caption{Rational coefficient approximations.  Each of the four coefficient columns is
scaled by $10^{10}$.}\label{tab:coefficients}\\
\toprule
$n$ & $10^{10}\widetilde c_{1,n}$ & $10^{10}\widetilde d_{1,n}$
 & $10^{10}\widetilde c_{2,n}$ & $10^{10}\widetilde d_{2,n}$\\
\midrule
\endfirsthead
\multicolumn{5}{c}{Table~\thetable\ (continued)}\\
\toprule
$n$ & $10^{10}\widetilde c_{1,n}$ & $10^{10}\widetilde d_{1,n}$
 & $10^{10}\widetilde c_{2,n}$ & $10^{10}\widetilde d_{2,n}$\\
\midrule
\endhead
\midrule
\multicolumn{5}{r}{Continued on the next page}\\
\endfoot
\bottomrule
\endlastfoot
1 & -670689128 & -6398628799 & -250018302 & 4107732082 \\
3 & 93524468 & 564897270 & -186215974 & -65204793 \\
4 & -64146977 & 121133118 & 82882809 & -327822296 \\
7 & -1358193 & -16329465 & 1979670 & 21637038 \\
9 & -2932910 & -940875 & 3636920 & 2145806 \\
12 & 3646952 & -34846939 & -6005382 & 34369042 \\
13 & 3183411 & 63030995 & -574123 & -75237042 \\
15 & 11390977 & 27049982 & -11312377 & -34985096 \\
16 & -5593246 & 13405071 & 5915710 & -10980954 \\
19 & 54905 & -1515007 & -120136 & 1498351 \\
21 & -210208 & -634307 & 177992 & 747946 \\
24 & 976077 & -804949 & -1037705 & 105643 \\
25 & -824472 & 7003719 & 1011732 & -6510015 \\
27 & 611891 & 4536095 & -428528 & -4440927 \\
28 & -472882 & -107613 & 399655 & 403751 \\
31 & 30744 & -112653 & -31773 & 90515 \\
33 & -4974 & -105568 & 1240 & 99081 \\
36 & 128313 & 246018 & -114876 & -267388 \\
37 & -176685 & 505187 & 166591 & -396209 \\
39 & -21638 & 472037 & 32128 & -402283 \\
40 & -18461 & -165597 & 10836 & 164974 \\
43 & 4565 & -2296 & -4050 & 547 \\
45 & 1477 & -10734 & -1524 & 8978 \\
48 & 10748 & 50050 & -8611 & -45048 \\
49 & -20506 & 7495 & 17354 & -266 \\
51 & -10361 & 29035 & 9641 & -21534 \\
52 & 2121 & -26523 & -2269 & 23159 \\
55 & 425 & 858 & -344 & -839 \\
57 & 302 & -632 & -266 & 463 \\
60 & 382 & 5658 & -231 & -4688 \\
61 & -1481 & -4320 & 1140 & 4098 \\
63 & -1494 & -352 & 1265 & 647 \\
64 & 558 & -2609 & -491 & 2114 \\
67 & 21 & 169 & -14 & -145 \\
69 & 34 & 13 & -28 & -18 \\
72 & -55 & 394 & 53 & -301 \\
73 & -21 & -805 & 5 & 685 \\
75 & -136 & -379 & 108 & 338 \\
76 & 70 & -144 & -58 & 104 \\
79 & -1 & 19 & 1 & -15 \\
81 & 2 & 9 & -2 & -8 \\
84 & -13 & 3 & 11 & 0 \\
85 & 13 & -87 & -11 & 70 \\
87 & -6 & -60 & 4 & 49 \\
88 & 6 & 5 & -4 & -5 \\
91 & 0 & 1 & 0 & -1 \\
93 & 0 & 1 & 0 & -1 \\
96 & -2 & -4 & 1 & 3 \\
97 & 2 & -6 & -2 & 4 \\
99 & 0 & -6 & 0 & 5 \\
100 & 0 & 2 & 0 & -2 \\

\end{longtable}
\endgroup

Write $J=\{1,3,4\}$ and $E=I_f\setminus J$.  Matrix rows and columns on a
node set always list all $c$ entries first and all $d$ entries second, with
nodes increasing within each group.  In particular the order on $J$ is
$(c_1,c_3,c_4,d_1,d_3,d_4)$.  The two rational matrices used to certify the
small inverses are
\begin{equation}\label{cert:Wplus}
 W_+=10^{-4}\begin{pmatrix}
 4159 & 4452 & -2003 & -10467 & -2421 & -1683 \\
2714 & 6573 & 1689 & 8258 & 104 & -335 \\
-4675 & 4785 & 7003 & -13100 & -1355 & -1163 \\
36853 & -37142 & 27514 & 105973 & 1100 & 2012 \\
57838 & -56941 & 36296 & 158623 & 29097 & 17293 \\
-4705 & 3589 & -3311 & -11932 & -2562 & 6919 \\

 \end{pmatrix},
\end{equation}
\begin{equation}\label{cert:Wminus}
 W_-=10^{-4}\begin{pmatrix}
 13265 & -1358 & -385 & -381 & 815 & -164 \\
616 & 10426 & 686 & -484 & 182 & 775 \\
-146 & 4 & 9842 & 215 & 267 & 179 \\
-2761 & 2483 & -958 & 5174 & -328 & 206 \\
-615 & -861 & 730 & -1388 & 5619 & -4370 \\
874 & 820 & 1148 & -431 & 905 & 11927 \\

 \end{pmatrix}.
\end{equation}
These matrices are exact inputs.  The $W$-based inverse argument below
uses their products and column sums; it does not require a computed
matrix inverse.  We describe the supplied program's additional inverse
checks in Section~\ref{cert:evaluation}.

\subsection{Moments and the integration rule}\label{cert:quadrature}

For a finite input $p=PR$ supported on nodes at most $100$, put
\begin{equation}\label{cert:moments}
 U_{p,l}(m,y)=\frac{y^l p^{(l)}(m)}{l!},\qquad
 X_{p,l}(m,y)=\frac{y^l K_p^{(l)}(m)}{l!}.
\end{equation}
The domain for the quadrature error estimate is stated explicitly in
Lemma~\ref{cert:moment-error}.

We use the first Fej\'er quadrature rule, whose cosine-sum weights are
described in \cite[Section~2, formula~(2.3)]{Waldvogel2006}. Our weights
are half the mass-two weights there, normalized to average over an interval;
for even $N=256$ the terminal cosine term in that formula vanishes.
We prove the exactness, positivity and error bound needed here.
Set $N=256$, and for $0\le v<N$ define
\begin{equation}\label{cert:rule}
\begin{aligned}
 \alpha_v&=\frac{v+1/2}{N},\qquad
 o_v=\frac1N\left(1-2\sum_{a=1}^{127}
             \frac{\cos(2a\pi\alpha_v)}{4a^2-1}\right),\\
 \tau_{jv}&=\frac{2j-1+\cos(\pi\alpha_v)}{12}
       \qquad(1\le j\le6).
\end{aligned}
\end{equation}
Integration over $(t_{j-1},t_j)$ is replaced by evaluation at $\tau_{jv}$
with weights $o_v/6$.  To check exactness, expand a polynomial in the
Chebyshev basis $T_k(\cos\theta)=\cos(k\theta)$.  Discrete cosine
orthogonality at $\theta_v=\pi\alpha_v$ recovers its coefficients through
degree $255$.  The integral averages of this basis on $[-1,1]$ are zero
at odd $k$ and $1/(1-k^2)$ at even $k$.  These are exactly the coefficients
in~\eqref{cert:rule}, so the rule is exact through degree $255$.
Moreover,
\[
 o_v\ge\frac1{256}\left(1-\sum_{a=1}^{127}\frac2{4a^2-1}\right)
     =\frac1{256\cdot255}>0,\qquad \sum_v o_v=1.
\]
Thus the six pieces together have total weight one.

For clarity, the complete density on piece $j$ used in this appendix is
\begin{equation}\label{cert:density}
 f_j(t)=2\pi\sum_n e^{-i\pi tn}
 \left\{-\pi c_n\sum_{l=j}^6(t_l-t)P_l e^{i\pi t_ln}
       +i d_n\sum_{l=j}^6 P_l e^{i\pi t_ln}\right\}.
\end{equation}
It includes node $0$ when that coefficient is part of the input.  The
folded atoms have masses $e_jCP_j$ at $t_j$, where $e_0=1$ and
$e_j=2$ for $j>0$.  The factors multiplying this density and these masses
for $U$ and $X$, respectively, are
\begin{equation}\label{cert:analytic-factors}
 e^{i\pi tm}\frac{(i\pi ty)^l}{l!},\qquad
 \lambda(t)e^{i\pi z(t)m}\frac{(i\pi z(t)y)^l}{l!}.
\end{equation}
The real part is taken after summation.

\begin{lemma}[Quadrature moment error]\label{cert:moment-error}
Let $p=PR$ have a real coefficient list supported in $\{0\}\cup I_f$,
with $\|(c,d)\|_1\le3$, and let $C\in\RR$ satisfy $|C|\le1$.
Let $m,y\in\RR$ and $l\in\{0,\ldots,30\}$ satisfy
$0\le m\le100$ and $|y|\le3/2$.
The quadrature rule~\eqref{cert:rule}, with exact evaluations at its nodes
and at the atoms, has error less than $10^{-22}$ in each moment
$U_{p,l}(m,y)$ and $X_{p,l}(m,y)$.
\end{lemma}

This domain contains all basis inputs and the two tabulated inputs used below.

\begin{proof}
Put $\rho=1/6$, $c_j=(2j-1)/12$, $a_j=c_j+ih$, and
$\Delta_j=|a_j|^2-\rho^2$.  These last quantities increase from
$\Delta_1=167/1200$ to $\Delta_6=389/400$, so the disk
$|t-c_j|\le\rho$ does not contain the pole $-ih$.  All factors in
\eqref{cert:density} and~\eqref{cert:analytic-factors} are analytic there.
The reciprocal-disk calculation is exact: for $\zeta=1/(t+ih)$,
\[
 |1-a_j\zeta|^2-\rho^2|\zeta|^2
 =\Delta_j\left|\zeta-\frac{\overline a_j}{\Delta_j}\right|^2
   -\frac{\rho^2}{\Delta_j}.
\]
Thus inversion maps the disk for $t+ih$ onto the disk with center
$(c_j-ih)/\Delta_j$ and radius $\rho/\Delta_j$.  In particular,
\[
 |\zeta|\le\frac{|a_j|+\rho}{\Delta_j}
       =\frac1{|a_j|-\rho}.
\]
The right side decreases with $c_j$, and $|a_1|=\sqrt{601}/60$.
Consequently the following are analytic bounds on every disk:
\[
\begin{aligned}
 |t|&\le\frac{13}{12}<1.1,\\
 |\lambda(t)|&\le\frac{120}{\sqrt3(\sqrt{601}-10)}<4.774<5,\\
 |z(t)|&\le\frac{80}{\sqrt{601}-10}+\frac25<5.912<6.2.
\end{aligned}
\]
For the two finer comparisons, it suffices to square the rational lower
bounds $1.732^2=2.999824<3$ and $24.515^2=600.985225<601$, and then use
\[
\begin{gathered}
 4.774(1.732)(24.515-10)=120.01826452>120,\\
 (5.912-.4)(24.515-10)=80.00668>80.
\end{gathered}
\]
The image disk for $z=-B\zeta-ih$ has center
$-B(c_j-ih)/\Delta_j-ih$ and radius $B\rho/\Delta_j$.  Its lowest
imaginary part is therefore $B(h-\rho)/\Delta_j-h$.  Since $h>\rho$,
this decreases with $\Delta_j$, giving
\begin{equation}\label{cert:disk-imaginary}
 \operatorname{Im}z(t)\ge
 B\frac{h-1/6}{c_j^2+h^2-(1/6)^2}-h
 \ge-\frac{1402}{17505}>-.081.
\end{equation}
Since $\sum_{l=1}^6|P_l|=10$ and $|t_l-t|\le13/12$ for $l\ge j$,
the analytic continuation of~\eqref{cert:density} on this disk is bounded
by
\[
 3\max\{2\pi^2\cdot10\cdot13/12,\,2\pi\cdot10\}
       e^{100\pi/6}<1000e^{100\pi/6}.
\]
Use $|a^l/l!|\le e^{|a|}$ for the divided powers.  Including all factors,
the direct and transformed integrands are bounded, respectively, by
\[
\begin{aligned}
 1000e^{200\pi/6+1.5\pi\cdot1.1}&<5.38\cdot10^{50},\\
 5000e^{100\pi/6+100\pi\cdot.081+1.5\pi\cdot6.2}
       &<1.51\cdot10^{50}.
\end{aligned}
\]
We may therefore use the common bound $M=10^{53}$.
The real interval has half-length $1/12$, half the disk radius.
Cauchy's coefficient estimate bounds the tail of the Taylor polynomial
of degree $255$ by $M\sum_{k=256}^{\infty}2^{-k}=M2^{-255}$.
Exactness, positivity and total mass one for the integral and quadrature
rule give discrepancy at most
\begin{equation}\label{cert:quadrature-error}
 2^{-254}10^{53}<3.455\cdot10^{-24}<10^{-22}.
\end{equation}
Atoms are evaluated separately and incur no quadrature error.
\end{proof}

\subsection{Interval evaluation}\label{cert:evaluation}

For a finite input $p$ in the real finite-input domain of
Lemma~\ref{cert:moment-error} and a node row $m\in I_f$, form the two real numbers
\begin{equation}\label{cert:pair}
 \mathcal S_m(p)=
 \left(-\frac{X_{p,0}(m,1)}{Q_m},\;
 \frac{D_mX_{p,0}(m,1)-X_{p,1}(m,1)}{Q_m}\right).
\end{equation}
With input $c_n=1$ or $d_n=1$ and all other data zero, these are the
entries of the corresponding column of $S$ at row node $m$.

Let $D=S_{J,J}$ and $U=S_{J,E}$.

For the Bernstein certificates, use exactly the triples $(i,m,y)$ and
orders $\nu$ specified in Lemma~\ref{signs:bernstein}; in particular $y$
is signed and no leftward interval at $m=1$ is used for index $1$.
Express its power coefficients in terms of the moments as
\begin{equation}\label{cert:power-from-moments}
 T_r=\sigma_i y^{-\nu}
 \left(U_{\widetilde p_i,r+\nu}(m,y)
             +X_{\widetilde p_{3-i},r+\nu}(m,y)\right).
\end{equation}
The signed scaling is exact.  Definition~\eqref{cert:moments} gives
$U_{p,l}(m,y)=y^lU_{p,l}(m,1)$ and the same identity for $X$.
Therefore, for positive or negative $y$,
\[
 T_r=\sigma_i y^r\left(U_{\widetilde p_i,r+\nu}(m,1)
                 +X_{\widetilde p_{3-i},r+\nu}(m,1)\right).
\]
This is the identity used when the implementation caches unscaled moments
at $y=1$; it does not discard the orientation of a half-gap.

Equations~\eqref{cert:rule}, \eqref{cert:density},
\eqref{cert:analytic-factors}, \eqref{cert:pair}, and
\eqref{cert:power-from-moments} already specify all
sums in the finite certificates below.  The following pseudocode records their
algebraic organization.  A measure is a list of pairs $(t,g)$ in the order
used by the supplied interval implementation: first the $6\cdot256$
quadrature pairs $(\tau_{jv},o_v f_j(\tau_{jv})/6)$, with
$j=1,\ldots,6$ outermost and $v=0,\ldots,255$ innermost, and then the seven
atom pairs $(t_j,e_jCP_j)$ for $j=0,\ldots,6$.  The atom masses are zero for
a basis input.  The pseudocode is not a byte-for-byte execution trace.
The supplied source fixes the detailed arithmetic and summation order for
an actual interval run, including the factoring and caching described below.

\noindent
\begin{minipage}{\linewidth}
\begin{lstlisting}[basicstyle=\ttfamily\footnotesize]
Moments(measure, m, y, L):
    U[0..L] = 0; X[0..L] = 0
    for (t,g) in measure:
        lam = i/(b*(t+i*h))
        z = -B/(t+i*h)-i*h
        u = g*exp(i*pi*t*m)
        v = g*lam*exp(i*pi*z*m)
        for r = 0..L:
            U[r] += real(u); X[r] += real(v)
            u *= (i*pi*t*y)/(r+1)
            v *= (i*pi*z*y)/(r+1)
    return U,X

Pair(measure,m):
    U,X = Moments(measure,m,1,1)
    return -X[0]/Q[m], (D[m]*X[0]-X[1])/Q[m]

for basis column c[n] or d[n], n in I_f:
    measure = Measure(basis column; all other data zero)
    column = stack Pair(measure,m), m in J
    if n in J:
        also form rows m in E intersect [7,16]
    test column sums of W[eta] and the defect/product matrices

for i in {1,2}, m in I_f:
    c,d = Pair(Measure(full table input 3-i),m)
    test abs(table_c[i,m]-c)
    test abs(table_d[i,m]+indicator(i==1 and m==1)-d)

for (i,m,y,nu) in the prescribed half-gap cases:
    U,X = Moments for the two full table inputs, through L=30
    T[r] = sigma[i]*(U[i,r+nu]+X[3-i,r+nu])/y**nu
    for k = 0..28:
        Bk = T[0]; a = 1
        for r = 1..k:
            a *= (k+1-r)/(29-r)
            Bk += a*T[r]
        test the appropriate lower bound for Bk
\end{lstlisting}
\end{minipage}

Here \texttt{stack} means all first entries followed by all second entries.
The recurrence for $a$ equals $\binom{k}{r}/\binom{28}{r}$.
\texttt{Measure} denotes this finite list and introduces no additional
approximation.  The fixed data at zero are used only for the two full
table inputs.  Reusing a measure means reusing its whole finite list.

For the accompanying certified evaluation, every real or complex quantity
in these sums is enclosed using interval arithmetic.  The pair calculation
forms the moments at $y=1$.  For the Bernstein calculation, the
implementation also forms and caches the unscaled moments at $y=1$, then
uses the signed $y^r$ identity following~\eqref{cert:power-from-moments}.
Each unscaled real moment enclosure is enlarged by
$[-10^{-22},10^{-22}]$, as justified by Lemma~\ref{cert:moment-error}.
Consequently the analytic-remainder contribution to the enclosure radius
of a Bernstein coefficient is bounded by
\[
\begin{aligned}
 2\cdot10^{-22}\sum_{r=0}^k
   \frac{\binom{k}{r}}{\binom{28}{r}}|y|^r
 &\le2\cdot10^{-22}\sum_{r=0}^{28}(3/2)^r\\
 &=4\cdot10^{-22}\bigl((3/2)^{29}-1\bigr)<5.2\cdot10^{-17}.
\end{aligned}
\]
The last inequality is an exact rational comparison.  It is not correct
to regard $10^{-22}$ as an unamplified per-Bernstein error.
Arb rounding and dependency effects can further widen the balls; all of
them are propagated by the subsequent products, absolute values, column
sums, and Bernstein combinations.  An upper bound is accepted only when
the complete enclosure lies strictly below the asserted rational bound,
and a lower bound only when it lies strictly above.  This procedure uses
enclosures of the exact integrals, not a quadrature error inferred from
observed numerical agreement.

The ball-arithmetic inclusion methodology is described in
\cite{Johansson2017}.  The supplied implementation uses Arb through
python-flint at $192$ bits of precision; the citation does not by itself
certify this local environment, its call domains, or a run.
The entry point \path{verification/verify_all.py} runs two verifiers in
the same directory, \path{interval_certificate.py} and
\path{interval_bounds.py}.  Their output records, \path{finite.json} and
\path{scalar.json}, and the aggregate record \path{summary.json} are
placed in its \path{results/} subdirectory.
The first verifier checks the exact bounds in~\eqref{cert:matrix-table},
\eqref{cert:exterior-block}, \eqref{cert:residual-bounds}, and
\eqref{cert:bernstein-table}, as well as all 2436 individual Bernstein
inequalities in Lemma~\ref{signs:bernstein}.  It also asserts
four finite brackets: the two defect brackets, the exterior bracket, and
the bracket for $(2,16,-1/2,28)$.  All its comparisons and extrema use
complete enclosures.

The mathematical $W$ argument in the next subsection uses only products
and column sums.  The supplied program also calls a certified inverse
for $I-\eta D$ and tests its inverse and cross-product norms as
additional checks; these calls are not needed for the $W$ deduction.

\subsection{Matrix and residual certificates}\label{cert:matrix-bounds}

The following are strict upper bounds
for the indicated column norms.  The last three columns use only
columns whose node lies in the stated range.  They refer to the exact
operator entries in~\eqref{cert:pair}.
\begin{equation}\label{cert:matrix-table}
\begin{array}{c|cc|ccc}
 \eta&
 \|W_\eta\|&
 \|I-W_\eta(I-\eta D)\|&
 \|W_\eta U_{7:9}\|&
 \|W_\eta U_{12:16}\|&
 \|W_\eta U_{19:100}\|\\ \hline
 +1&30.84&.001&3.54&.44&.25\\
 -1&1.83&.001&.73&.32&.17
\end{array}
\end{equation}
The unmultiplied exterior block also satisfies
\begin{equation}\label{cert:exterior-block}
 \|S_{E\cap[7,16],J}\|<.088.
\end{equation}
The certified enclosure calculation described in
Section~\ref{cert:evaluation} proves these inequalities.
For example, the two defect norms are enclosed respectively in
\[
\begin{gathered}
 [.00025422624910652544,\ .00025422624910652546],\\
 [.00042288165102553355,\ .00042288165102553357],
\end{gathered}
\]
and the exterior block norm is enclosed in
$[.08688748182289641591,\ .08688748182289641592]$.
These intervals are stated with rational endpoints rounded outwards.

The three printed product ranges partition $E$, so $\|W_\eta U\|$ is the
maximum of the three grouped product norms.  The exact bounds in
\eqref{cert:matrix-table} and~\eqref{cert:exterior-block} give the primary
matrix conclusion directly.  Put $A_\eta=I-\eta D$.
Since $\|I-W_\eta A_\eta\|<.001$, the finite square
matrix $W_\eta A_\eta$ is invertible by a geometric series, hence so
is $A_\eta$, and
\[
 A_\eta^{-1}=(W_\eta A_\eta)^{-1}W_\eta.
\]
It follows from~\eqref{cert:matrix-table} that
\[
 \|A_\eta^{-1}\|<\frac{30.84}{1-.001}<32,\qquad
 \|A_\eta^{-1}U\|<\frac{3.54}{1-.001}<3.7.
\]
Equation~\eqref{cert:exterior-block} gives the remaining bound $<.09$
in Lemma~\ref{interpolation:finite}.

For the two full tabulated inputs, including~\eqref{cert:fixed-inputs},
the residual certificate is
\begin{equation}\label{cert:residual-bounds}
\begin{aligned}
 \left|\widetilde c_{i,m}+\frac{X_{\widetilde p_{3-i},0}(m,1)}{Q_m}\right|
      &<3\cdot10^{-10},\\
 \left|\widetilde d_{i,m}+\mathbf1_{\{i=1,\ m=1\}}
       +\frac{X_{\widetilde p_{3-i},1}(m,1)
                  -D_mX_{\widetilde p_{3-i},0}(m,1)}{Q_m}\right|
      &<3\cdot10^{-10}
       \qquad(i=1,2,\ m\in I_f).
\end{aligned}
\end{equation}
The primary evaluation in Section~\ref{cert:evaluation} proves
\eqref{cert:residual-bounds} for these exact residuals.  Since
$3\cdot10^{-10}<10^{-9}$, these inequalities, the preceding primary
matrix conclusion, and~\eqref{cert:table-norms} prove every finite
assertion of Lemma~\ref{interpolation:finite}.

\subsection{Bernstein coefficient certificates}\label{cert:bernstein-bounds}

\begin{samepage}
Let $B_{i,m,y,k}$ denote the sum $B_k$ in
\eqref{signs:bernstein-bound}.  Each entry below is a strict lower bound
for every such sum in its row, for all $0\le k\le28$ and both half-gap
directions when applicable.
\begin{equation}\label{cert:bernstein-table}
\begin{array}{c|cc}
 m& i=1&i=2\\ \hline
 0&\text{---}&.762\\
 1&.314&.191\\
 3,4&.028&.024\\
 7,9&.114&.112\\
 12,13,15,16&.0105&.0094\\
 19,21&.123&.113\\
 24,25,27,28&.0111&.0113\\
 31,33&.121&.126\\
 36,37,39,40&.0107&.0115
\end{array}
\end{equation}
\end{samepage}
These bounds concern the exact derivatives of the tabulated functions,
not derivatives rounded to satisfy interpolation identities.
For example, the coefficient at $(i,m,y,k)=(2,16,-1/2,28)$ is enclosed in
\[
 [.00951875955680,\ .00951875955684].
\]
The primary evaluation in Section~\ref{cert:evaluation} proves the grouped
bounds~\eqref{cert:bernstein-table} for the exact derivatives of the tabulated
functions.  Each printed lower bound is stronger than its corresponding
threshold in~\eqref{signs:bernstein-bound}.  Therefore these exact grouped
bounds imply Lemma~\ref{signs:bernstein}.

\subsection{Scalar envelope and sign bounds}\label{cert:scalar-bounds}

The remaining finite substitutions involve elementary real expressions.
For the endpoint and density envelopes in
\eqref{interpolation:piece-envelopes}, it suffices to test the finitely
many endpoints and residues specified there.  The endpoint formulas
for $|\lambda|$, $|z|$, and $\operatorname{Im}z$ use rational numbers and
square roots.  The density expressions additionally use sine and cosine
at rational multiples of $\pi$.  The endpoint tests are certified by
the same elementary enclosures used above; their monotonicity and
convexity extensions to whole intervals were proved in
Section~\ref{sec:interpolation}.

For the row envelope in~\eqref{interpolation:row-envelope}, certified
outward bounds are
\begin{equation}\label{cert:row-substitutions}
\begin{array}{c|c}
 \text{row cut }x&\text{upper bound for the finite envelope sum}\\ \hline
 0&27.618723\\
 4&.269094\\
 16&.014122\\
 100&4.617010\cdot10^{-10}
\end{array}
\end{equation}
The atom-column analogue at cut $100$ is less than
$3.131186\cdot10^{-8}$.  These bounds imply all the row-tail constants
used in Section~\ref{sec:interpolation}.

For the sign proof, the following outward bounds record the relevant
substitutions.  The expressions themselves, and the reason they bound
whole intervals or tails, are given at the cited equations.
\begin{center}
\begin{tabular}{p{.65\linewidth}r}
\toprule
Expression & Certified bound\\
\midrule
\eqref{signs:rational-lower-sum}, index $1$ & $>.01206834$\\
\eqref{signs:rational-lower-sum}, index $2$ & $>.01656468$\\
$\sum_{j=1}^6 M_j(\pi t_j)^2/6$ & $<76.043$\\
Right side of \eqref{signs:tail-second-derivative}
before its final comparison & $<.005303$\\
\eqref{signs:perturbation-envelope} divided by $\nu!$,
$(s_*,\nu)=(0,0)$ & $<.003189$\\
Same expression, $(s_*,\nu)=(1/2,1)$ & $<.003736$\\
Same expression, $(s_*,\nu)=(1/2,2)$ & $<.009146$\\
Same expression, $(s_*,\nu)=(2,2)$ & $<.001796$\\
$2.29\sum_{j=1}^6M_j/6+.017\sum_{j=0}^6m_j^*$ & $<72.669$\\
\bottomrule
\end{tabular}
\end{center}
In the fourth row the expression means the explicit sum with
$.00005\cdot77$, not the derivative supremum to its left.
For the five cases of~\eqref{signs:taylor-representatives}, the
corresponding bounds for $\mathcal E(m,Y,\nu)$ in
\eqref{signs:taylor-bound} are
\begin{equation}\label{cert:taylor-substitutions}
\begin{array}{c|c}
 (m,Y,\nu)&\text{upper bound}\\ \hline
 (0,1/2,0)&1.578\cdot10^{-10}\\
 (1,1,1)&2.007\cdot10^{-5}\\
 (1,1,2)&5.735\cdot10^{-6}\\
 (3,1,2)&9.581\cdot10^{-11}\\
 (4,3/2,2)&6.133\cdot10^{-7}
\end{array}
\end{equation}
The first two rational-part bounds can be obtained by exact fraction
addition alone.  All the other displayed scalar bounds follow by
substitution and outward enclosure of the explicit exponential,
trigonometric and rational expressions.  The scalar verifier records
22 named strict comparisons: the four row substitutions, the atom tail,
the nine sign-table rows, the five Taylor representatives, and the three
disk-integrand and quadrature comparisons.  It also checks the endpoint
and residue envelopes, the coarser row-tail comparisons, the finite-inverse
substitution giving a bound below $90$,
and the distance substitution using $91$ and $2540$.  The derivation of
these whole-space coefficients from the Schur estimates remains the
analytic argument in Section~\ref{sec:interpolation}.
Its disk assertions use the coarser bounds $5$ and $6.2$ for
$|\lambda|$ and $|z|$; the finer $4.774$ and $5.912$ bounds were proved
analytically above and are not literal scalar-program assertions.
At the six gap midpoints the program evaluates the product defining $P$
and checks that each quotient is greater than $.68$; it does not assert
the symbolic radical identities for those values.  No sampling of an
unknown function between interpolation nodes is involved.  The
logarithmic concavity argument and exact midpoint values in
Lemma~\ref{signs:P-lower-bound} supply the global lower bound for $P$.
